% Transcription — fonds Alexandre Grothendieck, shelfmark 113, batch 1 (pages 1-12).
% Established from the facsimile archives/batches/113/batch-01.pdf (a mirror of
% https://grothendieck.umontpellier.fr/113.pdf), per the skill
% transcribe-grothendieck. The folder is twelve pages and this is its only
% batch.
%
% Eight of the twelve leaves are transcribed. Page 1 is an otherwise blank
% computer listing carrying two things in his hand: the single word
% « Abelianization » in the top right corner, and the number 377 in the top
% left. It is a cover, so it takes no \page marker and its word becomes the
% first section title; it is also where the inventory's title comes from, which
% is therefore his and not the archivists'. Page 2 is the other face of that
% leaf and is written landscape: it carries mathematics, and it is transcribed
% at its rank, but the question it treats is not the one the rest of the folder
% treats — it is about semi-simplicial MM-fibrés, submersions and an integral
% over a small category I, and nothing on it mentions abelianisation or
% k-additive categories. Pages 3, 5, 7, 9 and 11 are the run proper, in ink,
% carrying his own pagination — a circled 1 to 5 in the top right corner — so
% that run is continuous and complete. Page 12 is the reverse of page 11: a
% listing of the same run in whose two free margins he wrote, and it continues
% page 11 directly, so it is transcribed at its rank. Pages 4, 6, 8 and 10 are
% listings of that run with nothing of his on them and are skipped; the gaps in
% the page numbers are the record.
%
% The run of pages 3 to 12 is one construction taken in five steps, under a
% heading of his own boxed at the head of page 3: « (x) dans Cat », the tensor
% product of categories. Page 3 sets up P^k = Hom_Z(P°, Ab) for a small
% k-additive P, identifies Kar P with the projectives of P^k, and asks which
% abelian categories are of the form P^k — the answer being the abelian topoi
% with enough projectives, with a 2-functor (Add_k) --> (Top ab alg_k) written
% out and its morphisms described as the f_! commuting with colimits. Page 5
% carries the adjunction Hom_{k!}(P^k, M) ~ Hom_k(P, M) over into the two
% operations of P and Q = P° on P^k_M, and then defines P (x)_k Q by its
% 2-universal property Hom_k(P (x)_k Q, M) ~ Bil_k(P, Q; M), with the explicit
% construction by finite families and the two facts worth keeping. Page 7
% embeds P (x)_k Q in Bil_k(P°, Q°; Ab_k) by a (x)_k b |--> Hom(-,a) (x) Hom(-,b),
% and gives the first two applications: P^k_M and its dual are again
% k-abelian topoi, both computed as (P (x)_k Q)^k; and (P (x)_k Q)^k plays the
% role of the tensor product of the topoi P^k and Q^k. Page 9 proves that by
% currying, and opens the third application, the functor P^k (x)_k M -->
% Hom_k(P°, M) = P^k_M, shown fully faithful. Page 11 gives the two formulas
% that make L (x)_k x manageable — its convolution with L' and its Hom out of
% it — with the corollary that it is projective when both factors are, and
% closes on a question about P^k_M having enough projectives. Page 12 carries
% that question further: the derived F (x)^L_k L' by projective resolutions,
% then a list of four hypotheses on a k-additive abelian N (the fourth naming
% the projectives « k-plats ») and, in the other margin, the two propositions
% those hypotheses are meant to yield.
%
% Notation, fixed once for the whole batch. He draws two capital P's and two
% capital Q's: a plain one for the small k-additive category and a looped,
% script one for the abelian topos; the script pair is set here as \mathcal{P}
% and \mathcal{Q}, and \mathcal{P}_M abbreviates P^k_M. He abbreviates
% « stable » as « st.bl. » — read as « stable » throughout, which is the
% hypothesis every one of those isomorphisms needs. Duality is a small v set
% high to the right, distinct from his degree sign, and is set here as an
% exponent \vee. The sign « = » stands for the word « même » on pages 7 and 11,
% as it does in folder 116. On page 12 the convolution that page 11 writes *_k
% is drawn as an asterisk inside a circle, set here as \circledast; the circled
% star of page 9 is an equation label and not that operation.
%
% Readings flagged and not settled: the two letters before « Proj » on page 3,
% read « Ul »; the struck word or two after « Pour » on the same page; the word
% opening the second family on page 5, and the noun after « d'une » on page 7;
% the underlined abbreviation closing page 9, read « Can. » for « canonique »;
% and, on page 2, the word qualifying f in the factorisation g f, read
% « submersion » on the strength of the label « subm » that runs all over that
% leaf. Page 2's central web of M(...)'s exists twice, a first version cancelled
% in black and in red ink immediately to the left of the version transcribed
% here; the cancelled one is not given.
%
% The versos date the folder: the listings of pages 2 and 12 are stamped
% 16 SEP 82, and page 12's footer announces the machine's shutdown for Friday
% 17 September. That is presumably how the archivists arrived at
% « [à partir de 1982] », though nothing says so.
%
% Pass: Opus 5 (claude-opus-5), 2026-09-13 — first pass, unchecked against the pages by a human.
\documentclass[11pt,a4paper]{article}
% Le préambule commun à toutes les transcriptions du fonds Grothendieck.
%
% Il tient en une idée : **l'incertitude se compose, elle ne se résout pas.**
% Une transcription qui lisse un mot douteux a détruit la seule chose pour
% laquelle on la faisait — un lecteur doit pouvoir savoir, à chaque endroit,
% ce qui était sur le feuillet et ce qui a été ajouté. Les six macros qui
% suivent sont donc l'appareil critique, et non de la décoration.
%
% Les mêmes commandes sont reconnues par `scripts/render.mjs`, qui produit la
% vue de lecture du site. Une macro ajoutée ici doit l'être aussi là-bas,
% faute de quoi le rendu échouera bruyamment — ce qui est voulu : un rendu
% silencieusement faux est pire qu'un rendu absent.

\ProvidesPackage{grothendieck}[2026/08/08 Transcriptions du fonds Grothendieck]

\usepackage{amsmath,amssymb,amsthm}
\usepackage{mathrsfs}
\usepackage{stmaryrd}
% Les diagrammes commutatifs sont partout dans ce fonds : c'est la langue
% même de la géométrie algébrique de Grothendieck, et une transcription qui
% les rendrait en prose ne transcrirait rien.
\usepackage{tikz-cd}
% Français par défaut — c'est la langue des feuillets — mais l'anglais est
% chargé aussi : la traduction et le résumé sont des documents anglais, et la
% typographie française (espaces avant les deux-points) y serait une faute.
% Ces documents basculent par \selectlanguage{english} en tête de corps.
\usepackage[english,french]{babel}
\usepackage{fontspec}
\usepackage{geometry}
\geometry{a4paper,margin=25mm,marginparwidth=28mm}
\usepackage{marginnote}
\usepackage{xcolor}
% ulem, not soul. `soul`'s \st and \ul are built on a character-by-character
% scan that XeLaTeX's font machinery breaks: the document fails with an
% undefined control sequence rather than degrading. ulem does the same two jobs
% and is Unicode-engine safe. `normalem` stops it hijacking \emph, which the
% transcriptions use for Grothendieck's own emphasis.
\usepackage[normalem]{ulem}
\usepackage{hyperref}
\hypersetup{colorlinks=true,linkcolor=black,urlcolor=black,pdfborder={0 0 0}}

% --- Métadonnées du lot -----------------------------------------------------
% Reprises telles quelles par le script de rendu, qui les affiche en tête de
% la vue de lecture. Elles fixent ce que le fichier prétend couvrir : sans
% elles, une transcription flotte sans point d'attache dans un fonds de seize
% mille feuillets.

\newcommand{\folder}[1]{\gdef\grothFolder{#1}}
\newcommand{\batch}[1]{\gdef\grothBatch{#1}}
\newcommand{\leaves}[2]{\gdef\grothFirst{#1}\gdef\grothLast{#2}}
\newcommand{\foldertitle}[1]{\gdef\grothFoldertitle{#1}}
\gdef\grothFolder{?}\gdef\grothBatch{?}\gdef\grothFirst{?}\gdef\grothLast{?}\gdef\grothFoldertitle{}

% --- Le résumé --------------------------------------------------------------
%
% Il ouvre la lecture modernisée, sous le titre « Résumé », et il est composé
% en petit. Ce n'est pas un ornement : c'est le seul passage du document qui
% ne soit pas des mathématiques, et un lecteur doit pouvoir le reconnaître
% avant de l'avoir lu — puis le sauter s'il connaît déjà le sujet, ou s'y
% arrêter s'il ne le connaît pas. Le filet qui le suit marque la reprise du
% corps.

\newenvironment{resume}
  {\small\setlength{\parskip}{0.45em}}
  {\par\medskip\hrule\medskip}

% --- Mots-clés ---------------------------------------------------------------
%
% En anglais, en clôture du résumé de la lecture modernisée : le vocabulaire
% moderne sous lequel chercher ce que les feuillets construisent. Le manifeste
% du site (`npm run manifest`) les extrait pour étiqueter la cote entière —
% cette ligne est la source unique des tags, il n'y a pas de fichier à côté.
\newcommand{\keywords}[1]{\par\smallskip\noindent
  {\footnotesize\textsc{Keywords} --- \textit{#1}}\par}

% --- L'appareil critique ----------------------------------------------------

% Le numéro de feuillet, en marge. C'est l'ancre qui permet de régler un
% désaccord sur une formule en regardant *un* feuillet plutôt qu'une chemise
% de six cents. Le site s'en sert aussi pour tourner les pages du fac-similé
% quand on fait défiler la transcription.
\newcommand{\leaf}[1]{%
  \marginnote{\footnotesize\color{gray!70}\textsf{#1}}%
  \label{leaf:#1}\ignorespaces}

% Illisible : on ne devine pas. Un mot inventé qui se lit comme les autres est
% le pire résultat possible.
\newcommand{\ill}{\textcolor{red!60!black}{[\,\ldots\,]}}

% Lecture douteuse : le mot est donné, et signalé comme incertain.
\newcommand{\uncertain}[1]{\uline{#1}}

% Ajout de l'éditeur — une lettre manquante, un mot sous-entendu.
\newcommand{\add}[1]{\textcolor{blue!50!black}{[#1]}}

% Biffé par Grothendieck lui-même. Ce qu'il a rayé fait partie du document :
% une rature dit souvent où la pensée a changé de direction.
\newcommand{\struck}[1]{\sout{#1}}

% Note du transcripteur, sur l'état matériel du feuillet.
\newcommand{\note}[1]{\par\smallskip{\small\color{gray!60!black}\textit{[#1]}}\par\smallskip}

% Note marginale de Grothendieck, distincte du corps du texte.
\newcommand{\marginal}[1]{\par\smallskip\noindent\hspace{1em}%
  {\small\color{gray!60!black}#1}\par\smallskip}

% --- Titre ------------------------------------------------------------------

\AtBeginDocument{%
  \begin{center}
    {\large\bfseries Fonds Alexandre Grothendieck --- cote n\textsuperscript{o}~\grothFolder}\\[2pt]
    {\itshape\grothFoldertitle}\\[4pt]
    {\small feuillets \grothFirst--\grothLast\ (lot \grothBatch)}\\[2pt]
    {\footnotesize\color{gray!60!black}Transcription établie d'après le fac-similé numérisé
     par l'Université de Montpellier.\\
     Les lectures incertaines sont soulignées, les illisibles notées [\,\ldots\,],
     les ajouts entre crochets.}
  \end{center}
  \vspace{1em}}

\endinput


\folder{113}
\batch{1}
\pages{1}{12}
\dating{[à partir de 1982]}
\watermark{Édition de démonstration}
\foldertitle{Abelianization : notes manuscrites (s.d.)}

\begin{document}

\section{Abelianization}

\note{ce mot est de lui : il est seul, en haut à droite d'un listage d'ordinateur
par ailleurs vierge, à la page 1. Le feuillet ne portant que cela, il ne reçoit
pas de marque de page ; c'est de là que vient le titre de l'inventaire, qui est
donc de lui et non des archivistes. Le coin supérieur gauche du même feuillet
porte, de sa main, le nombre 377}

\section{Au dos du feuillet de titre : systèmes et MM-fibrés semi-simpliciaux}

\note{ce titre est de l'édition. Le feuillet n'en porte pas, et la question
qu'il traite n'est pas celle des feuillets 3 à 12 : il n'y est question ni
d'abélianisation ni de catégories $k$-additives. Il est écrit en travers, le
feuillet tourné d'un quart de tour}

\page{2}

\begin{tikzcd}[column sep=small, row sep=small, nodes={font=\scriptsize}]
\alpha \arrow[r, two heads] \arrow[d, "\xi"'] & \beta \arrow[d, dashed] & \gamma \arrow[l, hook'] \arrow[d, "\eta"] \\
X \arrow[r, two heads] & Y & Z \arrow[l, hook']
\end{tikzcd}

\begin{tikzcd}[column sep=small, row sep=small, nodes={font=\scriptsize}]
\alpha_0 \arrow[r, two heads] \arrow[d, "\xi_0"'] & \alpha'_0 \arrow[d, dashed, "\xi'_0"'] & \alpha_1 \arrow[l, hook'] \arrow[r] \arrow[d, "\xi_1"'] & \alpha_2 \arrow[r] \arrow[d, "\xi_2"'] & \cdots \arrow[r] & \alpha_{n-1} \arrow[r, two heads] \arrow[d, "\xi_{n-1}"'] & \alpha'_{n-1} \arrow[d, dashed, "\xi'_{n-1}"'] & \alpha_n \arrow[l, hook'] \arrow[d, "\xi_n"] \\
X_0 \arrow[r, two heads] & X'_0 & X_1 \arrow[l, hook'] \arrow[r] & X_2 \arrow[r] & \cdots \arrow[r] & X_{n-1} \arrow[r, two heads] & X'_{n-1} & X_n \arrow[l, hook']
\end{tikzcd}

\note{le feuillet porte en outre, entre les deux lignes de ce diagramme, une
série de flèches obliques portant les étiquettes $\eta_0, \eta_1, \dots,
\eta_{n-1}$ ; l'encombrement ne permet pas d'en séparer les extrémités et elles
ne sont pas tracées ici. Une accolade coiffe la ligne du haut, portant $u$ ;
une autre passe sous celle du bas, portant $v$}

Les syst.\ $(\xi'_0, \xi_{n-1})$ tels que $\xi'_{n-1} u = v \xi'_0$

\[
f \longmapsto \struck{\ill{}} \quad \struck{multibundle} \quad \text{MM-fibré} \quad M_1
\]

$M \longrightarrow$ \struck{foncteur} des syst.\ $\xi$.

\[
M \longrightarrow M_1
\]

\note{cette flèche, comme presque toutes celles de la page, porte l'étiquette
soulignée « subm »}

\[
M(\eta_0, \xi_{n-1}) \longrightarrow M(\xi'_0, \xi_n)
\]

\[
M\bigl(\xi'_0, \ill{}\, \eta_{n-1}, \xi'_{n-1}\bigr) \qquad M(\eta_0, \eta_{n-1})
\]

\[
M(\eta_0, \eta_1, \dots, \eta_{n-1})
\]

Puis, à droite :

\[
M\bigl(\xi_0, \xi'_0, \xi_1 \cdots \eta_{n-1}, \xi_{n-1}, \xi_n\bigr)
\;\simeq\; M(\xi_0, \xi_1, \dots, \xi_n)
\]
\[
M\bigl(\xi'_0, \xi_1, \xi_2 \cdots \xi_{n-1}, \xi'_n\bigr)
\longrightarrow \prod M(\eta_i, \xi_{i+1}, \eta_{i+1})
\qquad \text{si } i \leqslant n-2
\]
\[
M\bigl(\xi'_0, \eta_0, \eta_1, \dots, \eta_{n-1}, \xi'_n\bigr)
\longrightarrow \prod M(\eta_i, \eta_{i+1})
\qquad \text{à } i \leqslant n-2
\]
\[
M(\eta_0, \eta_{n-1}) \longleftarrow M(\eta_0, \eta_1, \dots, \eta_{n-1})
\qquad \text{subm } (\uncertain{mieux})
\]

\note{les listes d'indices de ces quatre lignes sont écrites très vite et leurs
apostrophes ne se séparent pas toutes ; elles sont données telles qu'elles se
lisent. Les flèches verticales entre elles portent « subm », l'une d'elles
« exact ». Immédiatement à gauche de ce bloc s'en trouve une première version,
biffée en noir puis en rouge, qui n'est pas transcrite}

Pour que $\int_I X(i)$ ne fasse pas sortir des MM-fibrés semi-simpliciaux, il
suffit que

\begin{itemize}
  \item[a)] Nerf.$(I)$ asphériques i.e.\ $I$ asphériques
  \item[b)] Chaque $X(i)$ transforme op.\ bord $d_i : \Delta_{n-1} \to \Delta_n$
    en submersions,
  \item[c)] $\forall\, u \in \mathrm{Fl}(I)$, $u : i \to j$,
    $X(u) : X(i) \to X(j)$ se factorise en $g f$, \ill{} $f$
    \uncertain{submersion}, $g$ mono.
\end{itemize}

\note{le coin inférieur droit du feuillet reprend le premier diagramme, sous la
forme $\alpha \twoheadrightarrow \beta \hookleftarrow \gamma$ au-dessus de
$X \twoheadrightarrow Y \hookleftarrow Z$, avec $\eta_0$ et $\eta_1$ pour
flèches verticales extrêmes ; les flèches du milieu y sont biffées}

\section{$\otimes$ dans Cat}

\note{ce titre est de lui : il est encadré au crayon en haut du feuillet 3. Les
feuillets 3, 5, 7, 9 et 11 portent en outre sa pagination propre, un chiffre
cerclé de 1 à 5 dans le coin supérieur droit}

\page{3}

$P$ petite catégorie $k$-additive

\[
P^{k} \;\underset{\text{déf}}{=}\; \operatorname{Hom}_{\mathbb{Z}}(P^{\circ}, Ab)
\;\simeq\; \operatorname{Hom}_k(P^{\circ}, Ab_k)
\]

(ne dépend pas du choix de la str.\ $k$-additive)

$P \hookrightarrow P^{k}$ inclusion additive pl.\ fid., d'où

\[
\operatorname{Kar} P \longrightarrow P^{k} \quad \text{qui induit équivalence}
\]

\marginal{N.B.\ $P \to \operatorname{Kar} P$ induit
$(\operatorname{Kar} P)^{k} \overset{\sim}{\to} P^{k}$}

\[
\operatorname{Kar} P \;\approx\; \uncertain{\mathrm{Ul}}\operatorname{Proj}(P^{k})
\]

\note{les deux lettres qui précèdent « Proj » sont lues « Ul » et ne sont pas
identifiées ; elles reparaissent plus bas sur le même feuillet}

Soit \struck{$P$} $\mathcal{P}$ catégorie \add{$k$-additive}, pour qu'elle soit
\add{$\approx$ cat.} de la forme $P^{k}$, avec $P$ \struck{additive}
$k$-additive, il f.\ et s.\ que

\begin{itemize}
  \item[a)] $\mathcal{P}$ soit un topos abélien, i.e.\ cat.\ abélienne, stable
    par petites $\varinjlim$ filtrantes, celles-ci étant exactes, et admettant
    une petite sous-cat.\ génératrice
  \item[b)] $\mathcal{P}$ a assez de projectifs.
\end{itemize}

\struck{\ill{}} et alors $\mathcal{P}$ on a \ill{} une $k$-équivalence
$\mathcal{P} \simeq P^{k}$, avec $P$ $k$-add.

\note{un mot ou deux sont biffés après « Pour », et « et alors » est écrit
au-dessus de la biffure ; ni l'un ni l'autre ne se lisent}

(\ill{} --- on peut prendre
$P = \uncertain{\mathrm{Ul}}\operatorname{Proj}(\mathcal{P})$).

\begin{tikzcd}[column sep=small, row sep=small, nodes={font=\scriptsize}]
(\mathrm{Add}_k) \arrow[r, "2\text{-ess. surj.}"] \arrow[d, "\operatorname{Kar}"'] & (\mathrm{Top\ ab}\ \ill{}_{\,k}) \\
(\mathrm{AddKar}_k) \arrow[ur, "3\text{-fid.}"']
\end{tikzcd}

\note{sous l'étiquette de la flèche horizontale, une seconde étiquette est
biffée, dont on lit le « 2 » initial. Le mot que porte le but, entre « ab » et
l'indice $k$, est repassé plusieurs fois et ne se lit pas ; la ligne qui suit
l'écrit « alg. »}

morphismes de topos $k$-abéliens alg.\ sont les
$f_{!} : \mathcal{P} \longrightarrow \mathcal{Q}$ commutant aux $\varinjlim$
telles que $f^{*}$ \struck{commut} ($=$ adjoint à droite) commutant aux
$\varinjlim$, i.e.\ admettant adjoint à droite $f_{*}$.

Soit $L \in \operatorname{Ob} P^{k}$, alors
$L \simeq \varinjlim_{P/L} a$

\[
\operatorname{Hom}_{k!}(P^{k}, M) \overset{\sim}{\longrightarrow}
\operatorname{Hom}_k(P, M)
\qquad \text{si $M$ $k$-additive \uncertain{stable} par $\varinjlim$}
\]

\note{il abrège « stable » en « st.bl. » ; c'est la lecture retenue partout
ci-dessous, et c'est l'hypothèse dont chacun de ces isomorphismes a besoin}

\struck{$\operatorname{Hom}_{\mathbb{Z}!}(P^{k}, Ab) \simeq
\operatorname{Hom}_{\mathbb{Z}}(P, Ab)$}

\[
\operatorname{Hom}_{k!}(P^{k\circ}, M) \overset{\sim}{\longrightarrow}
\operatorname{Hom}_k(P^{\circ}, M)
\qquad \text{si $M$ $k$-additive stable par $\varprojlim$}
\]

soit $Q = P^{\circ}$ :

\[
\operatorname{Hom}_{k!}(Q^{k\circ}, M) \simeq \operatorname{Hom}_k(P, M)
\]

\page{5}

\[
\underbrace{\operatorname{Hom}^{!}_{k}(Q^{k\circ}, M)}_{\textstyle \mathcal{Q}_M}
\;\simeq\;
\underbrace{\operatorname{Hom}_{k!}(P^{k}, M)}_{\textstyle \mathcal{P}_M}
\qquad \text{si $M$ $k$-add.\ stable pour les deux types de lim}
\]

Sur $\mathcal{P}_{M} \overset{\sim}{\longrightarrow}
\operatorname{Hom}_k(P^{\circ}, M)$ :

\begin{itemize}
  \item $\mathcal{Q} = \mathcal{P}^{\vee}$ opère par
    $(F, L') \mapsto F *_k L'$
  \item $\mathcal{P}$ opère par $(L, F) \mapsto \operatorname{Hom}_k(L, F)$
\end{itemize}

\note{l'exposant de $\mathcal{P}$ dans « $\mathcal{P}_M$ » est un $k$ biffé. La
marque de dualité est un petit $v$ placé haut et à droite, distinct de son signe
de degré $\circ$}

\[
\mathcal{P}^{k}_{M} \times Q^{k} \longrightarrow M,
\qquad
\mathcal{P}_M \times Q \longrightarrow M,
\qquad
P \times \mathcal{P}^{\dagger}_{M} \longrightarrow M,
\qquad
P^{k} \times P^{k}_{M}
\]

\note{ces quatre couples de facteurs sont écrits en colonne à droite des deux
lignes précédentes, comme un rappel des variances ; le signe porté par le
troisième ne se lit pas}

\uncertain{On} :
\[
(F *_k L')^{\circ} \;\simeq\; \struck{\ill{}}\ \operatorname{Hom}_k(L', F^{\circ}),
\qquad
\bigl(\operatorname{Hom}_k(L, F)\bigr)^{\circ} \;\simeq\; F^{\circ} * L
\]

\note{sous les membres de la première de ces deux formules il écrit $P, M$ et
$Q, M^{\circ}$, pour dire où chacun vit}

\note{un mot biffé, illisible, ouvre la seconde moitié du feuillet ; un trait
horizontal la sépare de ce qui précède}

Soient $P$, $Q$ $k$-additives, on définit $P \otimes_k Q$ et
$P \times Q \to P \otimes_k Q$ $k$-biadditive de façon 2-universelle :

\[
\operatorname{Hom}_k(P \otimes_k Q, M) \;\simeq\; \operatorname{Bil}_k(P, Q; M)
\qquad \text{tout $M$ $k$-additive}
\]

Construction explicite par familles finies d'objets de $P \times Q$ :

\[
\zeta = \bigoplus_{i \in I} a_i \otimes_k b_i ,
\qquad \ill{},\ \text{pour } \zeta' = \bigoplus_{j \in J} a'_j \otimes_k b'_j ,
\]
\[
\operatorname{Hom}(\zeta, \zeta') = \prod_{I \times J}
\operatorname{Hom}(a_i, a'_j) \otimes_k \operatorname{Hom}(b_i, b'_j)
\]

\textbf{À retenir}

\begin{itemize}
  \item[a)] Si $P$, $Q$ petites, $P \otimes_k Q$ est petite
  \item[b)] $\operatorname{Hom}_{P \otimes_k Q}(a \otimes_k b, a' \otimes_k b')
    \overset{\sim}{\longleftarrow}
    \operatorname{Hom}_P(a, a') \otimes_k \operatorname{Hom}_Q(b, b')$
\end{itemize}

On en conclut que le foncteur \add{can.} (pl.\ fid.\,!)

\note{« can » est inséré au-dessus de la ligne ; la phrase s'arrête là}

\page{7}

\[
P \otimes_k Q \;\hookrightarrow\; (P \otimes_k Q)^{k} \;\simeq\;
\struck{\operatorname{Bil}}\ \operatorname{Bil}_k(P^{\circ}, Q^{\circ}; (Ab_k))
\]

est donné par

\[
(*) \qquad a \otimes_k b \;\longmapsto\;
\bigl((x, y) \mapsto \operatorname{Hom}_P(x, a) \otimes_k
\operatorname{Hom}_Q(y, b)\bigr)
\]

ce qui donne une construction d'une \uncertain{variante} de $P \otimes_k Q$
comme la sous-cat.\ pleine de $\operatorname{Bil}_k(P^{\circ}, Q^{\circ}; Ab_k)$
formée des sommes finies de foncteurs de la forme $(*)$ (i.e.\ la cat.\ str.\
pleine correspondante)

\marginal{de l'existence de $P \otimes_k Q$, et du fait qu'elles petites si
$P$, $Q$ le sont}

\textbf{Application 1.} Soit $M \simeq Q^{k}$ un topos $k$-abélien alg.,
$M^{\vee} \simeq Q^{\circ k}$ « le » topos $k$-abélien alg.\ dual. Considérons
$P^{k}_{M}$ et $P^{\circ k}_{M^{\vee}}$ ($\simeq \mathcal{P}_M$ et
$(\mathcal{P}^{\circ})_{M^{\vee}}$), ce sont des deux topos $k$-abéliens alg.,
de $=$ \add{même}

\note{il écrit le signe « $=$ » pour le mot « même », ici et au feuillet 11}

En effet,
\[
\begin{aligned}
P^{k}_{M} &= \operatorname{Hom}_k\bigl(P^{\circ},
   \operatorname{Hom}_k(Q^{\circ}, Ab_k)\bigr) \\
 &\simeq \operatorname{Hom}_k(P^{\circ} \otimes_k Q^{\circ}, Ab_k) \\
 &\simeq \operatorname{Hom}_k\bigl((P \otimes_k Q)^{\circ}, Ab_k\bigr)
   = (P \otimes_k Q)^{k}
\end{aligned}
\]
\[
\begin{aligned}
P^{\circ k}_{M^{\vee}} &\simeq \operatorname{Hom}_k\bigl(
   (P^{\circ} \otimes_k Q^{\circ})^{\circ}, Ab_k\bigr) \\
 &\simeq \operatorname{Hom}_k(P \otimes_k Q, Ab_k) = (P \otimes_k Q)^{\circ k}
\end{aligned}
\]

\textbf{Appl.\ 2.}
\[
\operatorname{Bil}_{k!!}(P^{k}, Q^{k}; M) \;\simeq\;
\struck{\operatorname{Hom}}\ \operatorname{Hom}_{k!}\bigl((P \otimes_k Q)^{k}, M\bigr)
\]
\[
\bigl(\simeq \operatorname{Bil}_k(P, Q; M) \text{ si $M$ $k$-add.\ avec }
\varinjlim\bigr)
\]

\note{le $P$ du membre de droite porte un exposant $k$ biffé. Un point
d'interrogation est posé sous cet énoncé}

i.e.\ $(P \otimes_k Q)^{k}$ joue le rôle du prod.\ tens.\ des topos
$k$-abéliens $P^{k}$, $Q^{k}$, parmi les cat.\ $k$-additives stables par
$\varinjlim$

\page{9}

\emph{Dém.}
\[
\begin{aligned}
\operatorname{Bil}_{k!!}(P^{k}, Q^{k}; M)
 &\simeq \operatorname{Hom}_{k!}\bigl(P^{k},
    \underbrace{\operatorname{Hom}_{k!}(Q^{k}, M)}_{\textstyle
    \operatorname{Hom}_k(Q, M)}\bigr) \\
 &\simeq \operatorname{Hom}_k\bigl(P, \operatorname{Hom}_k(Q, M)\bigr) \\
 &\simeq \operatorname{Hom}_k(P \otimes_k Q, M)
    \overset{\sim}{\longleftarrow}
    \operatorname{Hom}_{k!}\bigl((P \otimes_k Q)^{k}, M\bigr)
\end{aligned}
\]

Dualement
\[
\operatorname{Bil}_{k!!}(P^{k\circ}, Q^{k\circ}; M) \;\simeq\;
\struck{\operatorname{Hom}}\ \operatorname{Hom}_{k!}\bigl(
(P \otimes_k Q)^{k\circ}, M\bigr)
\]

\note{ici encore le $P$ intérieur porte un exposant $k$ biffé}

\textbf{Appl.\ 3.} Soit $M$ $k$-additive stable par $\varinjlim$, on a

\[
(\circledast) \qquad
P^{k} \otimes_k M \longrightarrow \operatorname{Hom}_k(P^{\circ}, M) = P^{k}_{M}
\]
\[
L \otimes_k x \;\longmapsto\; \bigl(a \mapsto L(a) \otimes_k x\bigr)
= \widetilde{x} \circ L
\]

\note{le renvoi est une étoile inscrite dans un cercle, à distinguer de
l'astérisque nu du feuillet 7}

\[
P^{\circ} \overset{L}{\longrightarrow} Ab_k
\overset{\widetilde{x}}{\longrightarrow} M ,
\qquad \operatorname{Hom}_{k!}(Ab_k, M) \simeq M
\]

Ce foncteur $(\circledast)$ est pl.\ fidèle. En effet, \struck{comparons}.

\begin{tikzcd}[column sep=small, row sep=small, nodes={font=\scriptsize}]
M \arrow[r, "\text{pl. fid.}"] & M^{k} \\
\operatorname{Hom}_k(P^{\circ}, M) \arrow[r, "\text{pl. fid.}"] & \operatorname{Hom}_k(P^{\circ}, M^{k}) \arrow[d, "\text{pl. fid.}"] \\
& \operatorname{Hom}_k(P^{k\circ}, M^{k}) \\
& \operatorname{Bil}_k(P^{k} \times M^{\circ}, Ab_k)
\end{tikzcd}

\note{la troisième et la quatrième ligne sont reliées par un double trait
d'égalité, non par une flèche. À droite du diagramme il écrit : « attention, il
faut passer à un univers plus gros » ; et, contre la flèche verticale,
« NB $M^{k}$ st.\ $k$.\ par $\varprojlim$ ». Une première version de la
première ligne, biffée, portait
$\operatorname{Hom}_k(P^{\circ}, M) \to \operatorname{Hom}^{!}_{k}(P^{k\circ}, \ill{})$}

\[
P^{k} \otimes_k M \dashrightarrow \operatorname{Bil}_k(P^{k} \times M^{\circ}, Ab_k)
\qquad \text{pl.\ fid.\,?}
\]

\uncertain{Can.} connu pour être pl.\ fidèle

\note{trois jambages, lu « Can. » pour « canonique » ; « Com. » n'est pas exclu.
La flèche qui porte cette remarque est tracée en pointillé, et une seconde, en
pointillé également, monte vers le haut du diagramme}

\page{11}

On identifie\add{ra} parfois $L \otimes_k x$ à son image dans
$\operatorname{Hom}_k(P^{\circ}, M) = P^{k}_{M}$

\textbf{Ex.} $M = Q^{k}$, donc $P^{k}_{M} \simeq (P \otimes_k Q)^{k}$, donc
$(\circledast)$ devient

\[
P^{k} \otimes_k Q^{k} \longrightarrow (P \otimes_k Q)^{k}
\simeq \operatorname{Bil}_k(P, Q; Ab_k)
\qquad \text{pl.\ fid.}
\]
\[
F \otimes_k G \;\longmapsto\; \bigl((a,b) \mapsto F(a) \otimes_k G(b)\bigr)
\]

\note{sous « pl.\ fid. » une première étiquette, biffée, disait la même chose}

(On (cas dual, $=$ \add{même} si l'on \ill{} des variances $F$ : à \ill{} dans
$P$, $G$ dans $Q$, engendrant \struck{$\operatorname{Hom}$} $P^{k}_{M}$)

\textbf{Formule.} $L \in P^{k}$, $x \in M$ donc $L \otimes_k x \in P^{k}_{M}$ ;
$L' \in P^{\circ k}$.

\[
(L \otimes_k x) *_k L' \;\simeq\;
\underbrace{(L *_k L')}_{\textstyle \in Ab_k} \otimes_k x
\]

\note{il suffit de le vérifier pour $L \in \uncertain{P^{(a)}}$,
$L' \in \uncertain{P^{(b^{\circ})}}$ --- sous chacun des deux il pose un signe
$\|$ suivi de « $a$ » et de « $b^{\circ}$ », identifiant $L$ et $L'$ aux objets
dont ils sont les images ; les deux membres \ill{} $x$ fonctoriellement
$\simeq$ à \struck{\ill{}} $\operatorname{Hom}(\uncertain{b}, a) \otimes_k x$,
le premier argument de ce $\operatorname{Hom}$ étant repassé plusieurs fois}

Soit $F \in P^{k}_{M}$, alors
\[
\operatorname{Hom}_{P^{k}_{M}}(L \otimes_k x, F)
\;\simeq\; \operatorname{Hom}_M\bigl(x, \operatorname{Hom}_k(L, F)\bigr)
\;\simeq\; \operatorname{Hom}_{P^{k}}\bigl(L, \operatorname{Hom}_M(x, F)\bigr)
\]
où, par définition,
\[
\operatorname{Hom}_M(x, F) : \quad a \longmapsto
\operatorname{Hom}_M\bigl(x, F(a)\bigr)
\]

\marginal{si $L = a \in P$, \uncertain{on trouve}
$\operatorname{Hom}_M(x, F(a))$ pour les 2 derniers membres}

\note{cette remarque est écrite en biais, à environ 45°, le long du bord gauche
du feuillet}

\textbf{Cor.} Si $L \in \operatorname{Proj}(P^{k})$,
$x \in \operatorname{Proj}(M)$, alors
$L \otimes_k x \in \operatorname{Proj}(P^{k}_{M})$.

\textbf{Question.} Si $M$ $k$-add.\ stable par $\varinjlim$, et a assez de
projectifs, alors $P^{k}_{M}$ a assez de projectifs\add{ $F$}, et ceux-ci
annulent les $\operatorname{Tor}_i(L, F)$ ($L \in P^{k}$).

\note{le $F$ est inséré au-dessus de « projectifs », par un signe de renvoi}

\page{12}

De plus, on peut alors définir $F \circledast^{L}_{k} L'$ par résolutions proj.\
en $F$, comme $F \circledast_k L'$ --- un quasi-\uncertain{isom.} $L' \to L''$
définit un quasi-isom.\ $F \circledast_k L' \to F \circledast_k L''$.

\note{le signe de cette opération est un astérisque inscrit dans un cercle ;
c'est le produit que le feuillet 11 écrit $*_k$ sans le cercle}

Quand \ill{} et une \ill{} \ill{} pour les objets de $Q$ ainsi que pour
celles-ci \ill{}

\note{ces deux lignes sont écrites très pâle, entre deux traits ondulés qui
séparent ce bloc du suivant ; elles ne se lisent pas}

$N$ $k$-additive abélienne :

\begin{itemize}
  \item[a)] stable par petites sommes
  \item[b)] Une petite \ill{} de \ill{} \ill{} mono
  \item[c)] Tout objet de $N$ \ill{} quotient d'un projectif
  \item[d)] \add{Pour} tout objet projectif $x$ de $N$, \ill{} est
    « $k$-plat », i.e.\ (le foncteur
    \[
    Ab_k \longrightarrow N, \qquad M \longmapsto M \otimes_k a
    \]
    transforme mono en mono)
\end{itemize}

\note{« $k$-plat » est souligné et mis entre guillemets par lui ; c'est son nom
pour cette condition. Il écrit $a$ dans la formule là où la ligne parle de $x$,
et ce désaccord n'est pas réparé ici}

Dans l'autre marge du même feuillet :

Soit $N$ cat.\ $k$-additive abélienne stable par petites sommes.

\begin{itemize}
  \item[a)] Si dans $N$ tout objet est quotient d'un projectif, alors dans
    $P^{k}_{N}$ tout objet $F$ est quotient d'une somme
    $\Phi = \bigoplus_{a \in \operatorname{Ob} P} a \otimes_k x_a$, les $x_a$
    projectifs, et un tel $\Phi$ est projectif.
  \item[b)] Si de plus, \add{pour} tout objet projectif \struck{\ill{}} $x$ de
    $N$, le foncteur correspondant
    \[
    Ab_k \longrightarrow N, \qquad M \longmapsto M \otimes_k a
    \]
    est exact (i.e.\ transforme mono en mono), alors \struck{\ill{}} pour tout
    objet projectif $\Phi$ de $P^{k}_{N}$, le foncteur
    \[
    Q^{k} \longrightarrow N, \qquad L' \longmapsto \Phi *_k L'
    \]
    est exact --- et par suite $F \circledast^{L}_{k} L'$ peut être calculé ---
    redéfini par résolutions projectives $L'_i$ de $L'$\,! relativement
    projectives, et $F_i$ de $F$. \struck{De plus} \struck{$\ill{} \Phi
    \circledast_k L'$}
\end{itemize}

\note{ici encore la formule écrit $a$ là où la ligne parle de $x$. « relativement
projectives » est écrit sous la ligne, et la dernière ligne du feuillet est
barrée d'un long trait}

\end{document}
